\section*{The Appointed Formulary in Full}

Everything the 1962 \latin{Missale Romanum} prints for this Mass, in the order
it prints it: the Latin of the Vatican \latin{editio typica} exactly as the
controlling facsimile has it, each element under its own printed heading,
scriptural reference and marginal number, and beneath it the English of an
identified public-domain witness. This is the canonical full research edition;
a reader who came for the texts has them here.

\textbf{Formulary.} \latin{DOMINICA NONA post Pentecosten}, \latin{II
classis}, printed pp.~388--389, marginal nos.\ 1522--1531. Green follows the
general rubrics for Sundays after Pentecost and is not printed on those pages.
The formulary begins immediately after the Postcommunion of the Eighth Sunday
(no.\ 1521) and ends immediately before the heading \latin{DOMINICA DECIMA
post Pentecosten} (no.\ 1532). Nothing between those boundaries is omitted
here. It appoints no Tract, Sequence, second oration, blessing or ritual text,
and no second Collect, Secret or Postcommunion.

\textbf{The Latin.} Received liturgical text of the Roman Rite, older than the
edition that prints it. No public-domain conclusion is claimed for the 1962
typical edition itself; what permits the Latin here is 17 U.S.C.\ \S103(b),
under which a revised edition's copyright reaches only that edition's own
contribution and not the preexisting material it reprints. All ten appointed
elements of this formulary are printed in the Pustet Missal of Ratisbon, 1862,
which this repository tracks as public-domain text, and the terminal appendix
states the basis and its bounds in full. The project claims no rights in this text and has not composed,
translated, modernised, conflated or adapted any of it. Every form below was visually
collated at 300 and 600 dpi against the CMAA facsimile of the Vatican typical
edition, and independently against page images of the 1962 Benziger
\latin{editio iuxta typicam}. Where the two editions differ, the Vatican
typical reading is what appears here and the difference is tabulated in the
detailed commentary.

\textbf{The English.} Scriptural elements are the Douay--Rheims as revised by
Bishop Challoner (1749--1752) --- the English of the Clementine Vulgate the
missal prints --- in the Project Gutenberg e-text of that revision. The three
orations are the anonymous English of \work{The Roman Missal translated into
the English language for the use of the laity}, first revised edition
(Philadelphia: Eugene Cummiskey, 1861), from the three lines that book
supplies for this formulary, reproduced with its own nineteenth-century
spelling and its abbreviated conclusion ``Thro'.'' Both are public domain in
the United States. Neither is an official liturgical translation, and neither
is the text of the 1962 books: this is a study companion, not an altar book or
a hand missal. The project has translated nothing. Where the registered
English does not answer the missal's Latin --- and it fails to at four points
in this formulary --- the gap is marked at the element and left open.

\elementguard
\subsection*{1. Introit \cue{Int.}}

\begin{appointedlatin}{\latin{Ant. ad Introitum} --- Ps. 53, 6-7 --- marginal no.\ 1522}
Ecce Deus ádiuvat me, et Dóminus suscéptor est ánimæ meæ: avérte mala inimícis
meis, et in veritáte tua dispérde illos, protéctor meus, Dómine. \textnormal{Ps.
ibid., 3} Deus, in nómine tuo salvum me fac: et in virtúte tua líbera me.
℣.~Glória Patri.
\end{appointedlatin}

\begin{namedtranslation}{English: Douay--Rheims (Challoner), Ps. 53:6--7 and 53:3}
For behold God is my helper: and the Lord is the protector of my soul. Turn
back the evils upon my enemies; and cut them off in thy truth. \textit{Ps.}
Save me, O God, by thy name, and judge me in thy strength. ℣.~\latin{Glória
Patri.}
\end{namedtranslation}

\englishgap{Three clauses of the antiphon are not answered by that English, and
none of them is filled in here. The antiphon's closing vocative
\latin{protéctor meus, Dómine} has no counterpart anywhere in Psalm 53, so the
Douay has none either. The psalm verse sings \latin{líbera me} where the
Clementine, and therefore the Douay, reads \latin{judica me}: the English above
says ``judge me,'' the missal does not. And the Douay's opening ``For''
renders the \latin{enim} that the antiphon drops at v.~6. The doxology cue
\latin{Glória Patri} belongs to the Ordinary rather than to this formulary and
is a scriptureless liturgical formula; no English is supplied for it here.}

\elementguard
\subsection*{2. Collect \cue{Coll.}}

\begin{appointedlatin}{\latin{Oratio} --- marginal no.\ 1523}
Páteant aures misericórdiæ tuæ, Dómine, précibus supplicántium: et, ut
peténtibus desideráta concédas; fac eos, quæ tibi sunt plácita, postuláre. Per
Dóminum nostrum.
\end{appointedlatin}

\begin{namedtranslation}{English: Cummiskey (Philadelphia, 1861), Collect of this formulary, printed p.~414}
May the ears of thy mercy, O Lord, be open to the prayers of thy suppliants:
and, that they may succeed in their desires, make them ask those things that
are agreeable to thee. Thro'.
\end{namedtranslation}

\englishgap{Not a missing clause but a shifted subject, and it matters enough
to flag beside the text. \latin{Ut peténtibus desideráta concédas} makes God
the one who grants --- that \textit{thou} mayst grant to those who ask the
things they desire. The 1861 English makes the petitioners the ones who
succeed. The rendering is quoted because it is the registered public-domain
witness for this formulary; the argument built on the prayer's grammar is
conducted on the Latin, in the detailed commentary.}

\elementguard
\subsection*{3. Epistle \cue{Ep.}}

\begin{appointedlatin}{\latin{Léctio Epístolæ beáti Pauli Apóstoli ad Corínthios} --- 1 Cor. 10, 6-13 --- marginal no.\ 1524}
Fratres: Non simus concupiscéntes malórum, sicut et illi concupiérunt. Neque
idolólatræ efficiámini, sicut quidam ex ipsis: quemádmodum scriptum est: Sedit
pópulus manducáre et bíbere, et surrexérunt lúdere. Neque fornicémur, sicut
quidam ex ipsis fornicáti sunt, et cecidérunt una die vigínti tria mília.
Neque tentémus Christum, sicut quidam eórum tentavérunt, et a serpéntibus
periérunt. Neque murmuravéritis, sicut quidam eórum murmuravérunt, et
periérunt ab exterminatóre. Hæc autem ómnia in figúra contingébant illis:
scripta sunt autem ad correptiónem nostram, in quos fines sæculórum
devenérunt. Itaque qui se exístimat stare, vídeat ne cadat. Tentátio vos non
apprehéndat, nisi humána: fidélis autem Deus est, qui non patiétur vos tentári
supra id quod potéstis, sed fáciet étiam cum tentatióne provéntum, ut possítis
sustinére.
\end{appointedlatin}

\begin{namedtranslation}{English: Douay--Rheims (Challoner), 1 Cor. 10:6--13, from the point at which the lesson begins}
\dots\ that we should not covet evil things, as they also coveted. Neither
become ye idolaters, as some of them, as it is written: The people sat down to
eat and drink and rose up to play. Neither let us commit fornication, as some
of them committed fornication: and there fell in one day three and twenty
thousand. Neither let us tempt Christ, as some of them tempted and perished by
the serpents. Neither do you murmur, as some of them murmured and were
destroyed by the destroyer. Now all these things happened to them in figure:
and they are written for our correction, upon whom the ends of the world are
come. Wherefore, he that thinketh himself to stand, let him take heed lest he
fall. Let no temptation take hold on you, but such as is human. And God is
faithful, who will not suffer you to be tempted above that which you are able:
but will make also with temptation issue, that you may be able to bear it.
\end{namedtranslation}

\englishgap{The lesson opens mid-verse. Where the missal prints its own
liturgical address \latin{Fratres:}, the Clementine has \latin{Hæc autem in
figúra facta sunt nostri} and the Douay has ``Now these things were done in a
figure of us'' --- the words displaced above by the ellipsis. \latin{Fratres:}
is the missal's own and is not translated here.}

\elementguard
\subsection*{4. Gradual \cue{Grad.}}

\begin{appointedlatin}{\latin{Graduale} --- Ps. 8, 2 --- marginal no.\ 1525}
Dómine Dóminus noster, quam admirábile est nomen tuum in univérsa terra! ℣.
Quóniam eleváta est magnificéntia tua super cælos.
\end{appointedlatin}

\begin{namedtranslation}{English: Douay--Rheims (Challoner), Ps. 8:2}
O Lord, our Lord, how admirable is thy name in the whole earth! ℣.~For thy
magnificence is elevated above the heavens.
\end{namedtranslation}

\elementguard
\subsection*{5. Alleluia \cue{All.}}

\begin{appointedlatin}{Alleluia, no separate heading --- Ps. 58, 2 --- marginal no.\ 1526}
Allelúia, allelúia. ℣. \textnormal{Ps. 58, 2} Eripe me de inimícis meis, Deus meus:
et ab insurgéntibus in me líbera me. Allelúia.
\end{appointedlatin}

\begin{namedtranslation}{English: Douay--Rheims (Challoner), Ps. 58:2, between the chant's own acclamations}
Alleluia, alleluia. ℣.~Deliver me from my enemies, O my God; and defend me
from them that rise up against me. Alleluia.
\end{namedtranslation}

\elementguard
\subsection*{6. Gospel \cue{Gosp.}}

\begin{appointedlatin}{\latin{✠ Sequéntia sancti Evangélii secúndum Lucam} --- Luc. 19, 41-47 --- marginal no.\ 1527}
In illo témpore: Cum appropinquáret Iesus Ierúsalem, videns civitátem, flevit
super illam, dicens: Quia si cognovísses et tu, et quidem in hac die tua, quæ
ad pacem tibi, nunc autem abscóndita sunt ab óculis tuis. Quia vénient dies in
te: et circúmdabunt te inimíci tui vallo, et circúmdabunt te: et coangustábunt
te úndique: et ad terram prostérnent te, et fílios tuos, qui in te sunt, et non
relínquent in te lápidem super lápidem: eo quod non cognóveris tempus
visitatiónis tuæ. Et ingréssus in templum, cœpit eícere vendéntes in illo, et
eméntes, dicens illis: Scriptum est: Quia domus mea domus oratiónis est. Vos
autem fecístis illam spelúncam latrónum. Et erat docens cotídie in templo.
\end{appointedlatin}

\begin{namedtranslation}{English: Douay--Rheims (Challoner), Lk. 19:41--47a}
And when he drew near, seeing the city, he wept over it, saying: If thou also
hadst known, and that in this thy day, the things that are to thy peace: but
now they are hidden from thy eyes. For the days shall come upon thee: and thy
enemies shall cast a trench about thee and compass thee round and straiten
thee on every side, And beat thee flat to the ground, and thy children who are
in thee. And they shall not leave in thee a stone upon a stone: because thou
hast not known the time of thy visitation. And entering into the temple, he
began to cast out them that sold therein and them that bought. Saying to them:
It is written: My house is the house of prayer. But you have made it a den of
thieves. And he was teaching daily in the temple.
\end{namedtranslation}

\englishgap{The missal's \latin{In illo témpore: Cum appropinquáret Iesus
Ierúsalem} names a subject and a city that Luke does not name at v.~41; the
Clementine reads \latin{Et ut appropinquávit} and the Douay accordingly begins
``And when he drew near.'' The liturgical incipit is the missal's own and is
not translated here. The reading stops at v.~47a: the Douay's v.~47 continues
``And the chief priests and the scribes and the rulers of the people sought to
destroy him,'' and the missal cuts before it.}

\appointedrubric{Credo.}{Printed on its own line immediately after the Gospel.
A direction, not a text: the Creed itself is in the Ordinary.}

\elementguard
\subsection*{7. Offertory \cue{Off.}}

\begin{appointedlatin}{\latin{Antiphona ad Offertorium} --- Ps. 18, 9, 10, 11 et 12 --- marginal no.\ 1528}
Iustítiæ Dómini rectæ, lætificántes corda, et iudícia eius dulcióra super mel
et favum: nam et servus tuus custódit ea.
\end{appointedlatin}

\begin{namedtranslation}{English: Douay--Rheims (Challoner), Ps. 18:9--12 in full, with the four fragments the antiphon takes in bold}
\textbf{v.~9} \textbf{The justices of the Lord are right, rejoicing hearts:}
the commandment of the Lord is lightsome, enlightening the eyes.
\textbf{v.~10} The fear of the Lord is holy, enduring for ever and ever:
\textbf{the judgments of the Lord} are true, justified in themselves.
\textbf{v.~11} More to be desired than gold and many precious stones: and
\textbf{sweeter than honey and the honeycomb.} \textbf{v.~12} \textbf{For thy
servant keepeth them,} and in keeping them there is a great reward.
\end{namedtranslation}

\englishgap{This antiphon is not a quotation but a centonisation, which is why
the missal prints four verse numbers instead of a range: it takes v.~9a, the
subject \latin{iudicia} from v.~10c, the honey comparison from v.~11b, and
v.~12a. The four Douay verses are therefore printed whole with the fragments
marked, rather than stitched into a seamless English paragraph that the
project would itself have composed. Read in the antiphon's own order the
fragments say: the justices of the Lord are right, rejoicing hearts, and the
judgments of the Lord sweeter than honey and the honeycomb; for thy servant
keepeth them. One further difference is invisible in English: the missal
writes \latin{nam et} where the psalm has \latin{Etenim}, and the Douay reads
``For'' either way.}

\elementguard
\subsection*{8. Secret \cue{Sec.}}

\begin{appointedlatin}{\latin{Secreta} --- marginal no.\ 1529}
Concéde nobis, quæsumus, Dómine, hæc digne frequentáre mystéria: quia,
quóties huius hóstiæ commemorátio celebrátur, opus nostræ redemptiónis
exercétur. Per Dóminum nostrum Iesum Christum, Fílium tuum: Qui tecum vivit et
regnat in unitáte.
\end{appointedlatin}

\begin{namedtranslation}{English: Cummiskey (Philadelphia, 1861), Secret of this formulary, printed p.~415}
Grant us, O Lord, we beseech thee, frequently and worthily to celebrate these
mysteries: for as many times as this commemorative sacrifice is celebrated, so
often is the work of our redemption performed. Thro'.
\end{namedtranslation}

The typical edition prints the long conclusion here, breaking off at
\latin{in unitáte} for the celebrant to complete from the Ordinary; Benziger
prints only \latin{Per Dóminum}, and the 1861 hand missal only ``Thro'.'' The
1861 rendering splits the single verb \latin{frequentáre} and the adverb
\latin{digne} into ``frequently and worthily to celebrate.''

\appointedrubric{Præfatio de Ssma Trinitate.}{Printed immediately after the
Secret, with \latin{Ssma} under a tilde abbreviation for \latin{Sanctissima}:
the Preface of the Most Holy Trinity, which is in the Ordinary. Expanded here
from the same book's own abbreviation practice.}

\elementguard
\subsection*{9. Communion \cue{Comm.}}

\begin{appointedlatin}{\latin{Ant. ad Communionem} --- Io. 6, 57 --- marginal no.\ 1530}
Qui mandúcat meam carnem, et bibit meum sánguinem, in me manet, et ego in eo,
dicit Dóminus.
\end{appointedlatin}

\begin{namedtranslation}{English: Douay--Rheims (Challoner), Jn. 6:57}
He that eateth my flesh and drinketh my blood abideth in me: and I in him.
\end{namedtranslation}

\englishgap{The antiphon's closing \latin{dicit Dóminus} is the liturgy's own
attribution formula, not part of John 6:57, and therefore not in the Douay.
The antiphon also reads \latin{et ego in eo} where the Clementine has
\latin{et ego in illo}; the English is ``and I in him'' either way. The
reference \latin{Io. 6, 57} is the Clementine's verse number, which the
Douay--Rheims shares; most current editions number the same clause 6:56.}

\elementguard
\subsection*{10. Postcommunion \cue{Postcomm.}}

\begin{appointedlatin}{\latin{Postcommunio} --- marginal no.\ 1531}
Tui nobis, quæsumus, Dómine, commúnio sacraménti, et purificatiónem cónferat,
et tríbuat unitátem. Per Dóminum nostrum.
\end{appointedlatin}

\begin{namedtranslation}{English: Cummiskey (Philadelphia, 1861), Postcommunion of this formulary, printed p.~415}
May the participation of this thy sacrament, O Lord, we beseech thee, both
purify us, and unite us. Thro'.
\end{namedtranslation}

\clearpage
